\newpage
\chapter{この手引きについて}

\begin{enumerate}[label=\thechapter.\arabic*　,leftmargin=*]

  \item \textbf{この手引きの位置づけ}

  学園の文書は三層で構成され、学園憲章が原則を、学園規則が制度を、ガイドライン類が運用を定める。
  この手引きはガイドライン層に属し、教務主事が職務を行うときの指針を定める。

  この手引きは、学園規則第４章第１条第２項にいう教務主事の職務に関する具体的な手続に当たる。

  \item \textbf{学園規則第４章との関係}

  学園規則第４章は、教務主事の基本理念として、支援、対等、抑制、説明及び継承の五つの原則を定めている。
  この手引きは、これらの原則を日々の職務でどう実行するかを示すものであり、新たな原則を立てるものではない。

  この手引きは、学園規則に定めのない義務を学生に課さない。

  \item \textbf{この手引きに定めのない事項}

  この手引きは起こりうるすべての事態を網羅しない。
  そのため、定めのない事態に直面したときは、学園規則第４章の原則に立ち返って判断し、その手順は第３章に定める。

  文書の様式、連絡の文例その他の実務上の細目は、この手引きとは別に定める。

  \item \textbf{この手引きの改定}

  この手引きは運用の実際に合わせて随時改定するものとし、学園規則の改定手続を要しない。

  ただし、学園憲章又は学園規則に反する内容は、その効力を有しない。
  改定するときは、上位の文書との整合を確認する。

\end{enumerate}